%%%%%%%%%%%%%%%%%%%%%%%%%%%%%%%%%%%%%%%%%%%%%%%%%%%%%%%%%%%%%%%%%%%%%%%%
%                                                                      %
%     File: Conclusions/01_Achievements.tex                            %
%     Tex Master: ../Thesis.tex                                        %
%                                                                      %
%%%%%%%%%%%%%%%%%%%%%%%%%%%%%%%%%%%%%%%%%%%%%%%%%%%%%%%%%%%%%%%%%%%%%%%%

% ==============================================================================
\section{Achievements}
\label{sec:concl_summary}

% What did the work set out to do?
This work set out to operate a \gls{usv} and a \gls{uav} connected by a tether as a single system at sea, each vehicle within its own limits and both within the reach of the tether, under waves and wind. It pursued this through the four objectives of Section~\ref{sec:intro_problem_statement}, whose outcomes are reviewed in the order in which the work built on them.

% How well do the models describe the system?
The models were derived from the physics of each vehicle and validated against the simulator with its nominal parameters, without any fitting. Over the horizons of the controllers, both predict the motion of their vehicle with errors of a few percent of its range, small enough for the controllers to anticipate it. The tether, often treated in the literature as massless and taut or as an unknown disturbance, is described by a closed-form model of the force of a slack tether that accounts for its weight, which reduces the error of the force by more than two thirds relative to those simpler descriptions, at a cost low enough to be evaluated along the whole horizon.

% What do the controllers achieve?
On these models, each vehicle is driven by its own \gls{nmpc}, and the two exchange the trajectories they predict, so that each keeps the separation within the reach of the tether while anticipating where the other vehicle is going. Both solve their problems well within their sampling periods. On a trajectory that the \gls{uav} can follow, its controller tracks it with an error more than four times smaller than that of the PX4 position controller, since it anticipates the motion instead of reacting to the error. When a reference demands more than a vehicle can deliver, the controller gives up following it only as far as the limits require, and uses the full capability of the vehicle in doing so, so that the manoeuvres can remain simple and give where each vehicle should be rather than how to get there.

% Does the landing hold at sea?
The landing on the moving deck, the manoeuvre most exposed to the sea, holds under waves and wind. Since the descent is referenced to the height of the winch rather than to the heaving deck, the \gls{uav} touches down a few centimetres from the centre of the pad in every sea tested, and the waves change only the speed of the contact. Under wind, the disturbance observer keeps the \gls{uav} centred above the pad and the landing completes as without wind, whereas without it the \gls{uav} is never centred and the landing is aborted.

% Do the two vehicles work as one system?
Coupled by the tether, the two vehicles complete the three missions as a single system. In the inspection, the \gls{uav} surveys a tower from bottom to top while the \gls{usv} carries the tether around it from the water, which neither vehicle could do alone. In the patrol, the \gls{uav} keeps a fleeing target in view, waiting at the limit of the tether while the \gls{usv} closes the distance, and in the recovery, when the \gls{uav} cannot follow, the \gls{usv} moves no further than the tether allows and returns to collect it. Each mission ends with the \gls{uav} landing back on the moving deck, close to the centre of the pad.

% What do these results show as a whole?
Together, these results show that two cooperating \gls{nmpc} controllers, built on models of both vehicles and of the tether, can operate a tethered \gls{usv} and \gls{uav} as a single system at sea, in simulation. Through this cooperation, the tether becomes not only a constraint to be avoided, but a means to extend what each vehicle can do alone, as the inspection of the tower shows.
